% Ableitungsregeln – bank/ableitungsregeln/ (zusammenbau v0.4, Heft)
\einheitenkopf[t2]{Ableitungsregeln}
\setcounter{aufgabe}{11}

% Anlauf (ohne Prüfkennung)
\begin{aufgabe}{Potenz-, Faktor- und Summenregel – Anlauf}
\begin{teile}
\teil $f(x) = x^9$ – welche Regel reicht? \\ \kreuz{Potenzregel} \\ \kreuz{Faktorregel} \\ \kreuz{Konstantenregel}
\teil $f(x) = 2x^3 + 5x^2 - 4x + 9$ – $f'(x)$? $f'(x) =$ \leerfeld
\teil $f(x) = \frac{1}{6}x^3 + \frac{3}{2}x^2 - 2x$ – $f'(x)$? $f'(x) =$ \leerfeld
\end{teile}
\end{aufgabe}

\begin{aufgabe}{Potenz-, Faktor- und Summenregel – Prüfungsaufgaben}
\begin{teile}
\teil $f(x) = -3x^5 + 4x^3 - x$ – Ableitung $f'(x)$ und eine Stammfunktion $F(x)$? \hfill (Abitur 2019 GK) \\ $f'(x) =$ \leerfeld , $F(x) =$ \leerfeld
\teil $f(x) = 6x^2 - 2x + 5$ – Ableitung $f'(x)$ und eine Stammfunktion $F(x)$? \hfill (Abitur 2019 GK) \\ $f'(x) =$ \leerfeld , $F(x) =$ \leerfeld
\teil $f(x) = 8x^3 + 3x^2 - 4x$ – Ableitung $f'(x)$ und eine Stammfunktion $F(x)$? \hfill (Abitur 2019 GK) \\ $f'(x) =$ \leerfeld , $F(x) =$ \leerfeld
\teil $f(x) = \frac{1}{5}x^5 - 6x^3 + 81x$. Weise nach: $f'(x) = (x - 3)^2 \cdot (x + 3)^2$. \hfill (Abitur 2022 GK)
\teil $f(x) = \frac{1}{20}x^5 - \frac{8}{3}x^3 + 64x$. Weise nach: $f'(x) = \frac{1}{4} \cdot (x - 4)^2 \cdot (x + 4)^2$. \hfill (Abitur 2022 GK)
\teil $f(x) = \frac{2}{25}x^5 - \frac{20}{3}x^3 + 250x$. Weise nach: $f'(x) = \frac{2}{5} \cdot (x - 5)^2 \cdot (x + 5)^2$. \hfill (Abitur 2022 GK)
\end{teile}
\end{aufgabe}

% Anlauf (ohne Prüfkennung)
\begin{aufgabe}{Produktregel – Anlauf}
\begin{teile}
\teil $f(x) = x^3 \cdot e^x$ – welche Regeln? \\ \kreuz{nur Produktregel} \\ \kreuz{Produkt- und Kettenregel} \\ \kreuz{nur Kettenregel}
\teil $f(x) = x \cdot e^x$ – $f'(x)$, $e^x$ ausgeklammert? $f'(x) =$ \leerfeld
\end{teile}
\end{aufgabe}

\begin{aufgabe}{Produktregel – Prüfungsaufgaben}
\begin{teile}
\teil $f(x) = 3 \cdot (x^2 - 4) \cdot e^x$. Weise nach: $f'(x) = 3 \cdot (x^2 + 2x - 4) \cdot e^x$. \hfill (Abitur 2024 GK)
\teil Der Graph von $f$ mit $f(x) = (2x + 1) \cdot e^{-x}$ hat genau einen Extrempunkt. Weise nach: $f'(x) = (-2x + 1) \cdot e^{-x}$, und berechne die Koordinaten des Extrempunkts auf zwei Stellen gerundet. \hfill (Abitur 2020 GK) \\ Extrempunkt: ( \leerfeld | \leerfeld )
\teil Die Abbildung zeigt den Graphen der linearen Funktion $g$ und den Graphen von $f$ mit dem Tiefpunkt $T$. Für $h(x) = f(x) \cdot g(x)$: Steigung der Tangente an den Graphen von $h$ an der Stelle $x = 2$? \hfill (Abitur 2026 GK) \\

\begin{ksys}[xmin=-1,xmax=6,ymin=-2,ymax=6,ablesen] \gerade{-1}{5}{g} \parabel{1}{2}{-1}{f} \tiefpunkt{2}{-1}{T} \end{ksys}
$h'(2) =$ \leerfeld
\teil Die Abbildung zeigt den Graphen der linearen Funktion $g$ und den Graphen von $f$ mit dem Tiefpunkt $T$. Für $h(x) = f(x) \cdot g(x)$: Steigung der Tangente an den Graphen von $h$ an der Stelle $x = 1$? \hfill (Abitur 2026 GK) \\

\begin{ksys}[xmin=-2,xmax=5,ymin=-4,ymax=5,ablesen] \gerade{0.5}{1.5}{g} \parabel{1}{1}{-3}{f} \tiefpunkt{1}{-3}{T} \end{ksys}
$h'(1) =$ \leerfeld
\teil Die Abbildung zeigt den Graphen der linearen Funktion $g$ und den Graphen von $f$ mit dem Tiefpunkt $T$. Für $h(x) = f(x) \cdot g(x)$: Steigung der Tangente an den Graphen von $h$ an der Stelle $x = 1$? \hfill (Abitur 2026 GK) \\

\begin{ksys}[xmin=-1,xmax=4,ymin=-1,ymax=7,ablesen] \gerade{-2}{6}{g} \parabel{1}{1}{2}{f} \tiefpunkt{1}{2}{T} \end{ksys}
$h'(1) =$ \leerfeld
\end{teile}
\end{aufgabe}
